% Zdroj: PDF priklady-leto-2023.pdf, fyzická str. 6-7. Značka: otázka studijního zaměření. Zařazeno: I4.
% Stránkovací tabulka je oříznutá z PDF (str. 7) do fig/strankovaci-tabulka.png.
\section{Překlad adresy}
\szzotazka{léto 2023}{17}{Překlad adresy (stránkování)}
\noindent\textit{(Pozn.: v~rámci současné akreditace nešlo o~společnou otázku (q1--q4), ale o~otázku jiné specializace (specializační sekce, systémové zaměření, ne UI). Do společného okruhu I4 tedy formálně nepatří; tematicky sem ovšem spadá a~je výborná k~procvičení.)}\par

\noindent\textbf{Zadání.}\par
Předpokládejme hypotetický procesor s~paměťovou architekturou, kde virtuální i~fyzický
prostor používá 16-bitové adresování. Pro překlad adres se používá 1-úrovňové stránkování
s~64\,B stránkami. Záznamy ve stránkovací tabulce (které mají také 16 bitů) využívají
offsetové bity pro přidružené příznaky, nejméně významný bit je použit jako příznak
platnosti (tj. 1 = záznam je platný, 0 = záznam není platný).

Stránkovací tabulka, jejíž výpis následuje níže, se nachází na adrese \texttt{0x1d40}.
Výpis zobrazuje pro jednoduchost přímo 16-bitové položky stránkovací tabulky, tedy není
nutné řešit pořadí bytů v~reprezentaci čísel.

\begin{center}\includegraphics[width=\linewidth]{strankovaci-tabulka.png}\end{center}

Přeložte virtuální adresu \texttt{0x328f} (\texttt{0b0011001010001111} binárně) na fyzickou adresu:
\begin{enumerate}
\item Napište číslo stránky a~offset pro danou virtuální adresu.
\item Napište fyzickou adresu položky ve stránkovací tabulce, která odpovídá zadané virtuální adrese.
\item Napište výslednou fyzickou adresu, pokud existuje překlad. Jinak uveďte, že neexistuje.
\end{enumerate}

V~odpovědi důsledně uvádějte číselnou soustavu -- čísla bez prefixu jsou v~desítkové
soustavě, prefix \texttt{0x} označuje šestnáctkovou soustavu, prefix \texttt{0b} označuje
dvojkovou soustavu. (To, jakou soustavu použijete, se ovšem samo o~sobě nehodnotí.)

\medskip
\noindent\textbf{Náčrt řešení.}\par
\begin{enumerate}
\item Stránka má $64=2^{6}$\,B, takže offset jsou dolní \textbf{6} bitů a~číslo stránky horních \textbf{10} bitů:
  \[\texttt{0x328f}=\underbrace{\texttt{0b0011001010}}_{\text{číslo stránky}}\ \underbrace{\texttt{001111}}_{\text{offset}},\]
  číslo stránky $=\boxed{\texttt{0xCA}=202}$, offset $=\boxed{\texttt{0x0F}=15}$.
\item Položky tabulky jsou 16-bitové (2\,B), položka stránky $202$ je na adrese
  $\texttt{0x1d40}+202\cdot 2=\texttt{0x1d40}+\texttt{0x194}=\boxed{\texttt{0x1ED4}}$.
\item Na adrese \texttt{0x1ED4} je ve výpisu hodnota \texttt{0xFBF7}. Nejnižší bit je $1$, položka je
  \emph{platná}. Číslo rámce tvoří horních $10$ bitů (dolních $6$ jsou příznaky): $\texttt{0xFBF7}\,\&\,\texttt{0xFFC0}=\texttt{0xFBC0}$.
  Fyzická adresa vznikne spojením rámce s~offsetem: $\texttt{0xFBC0}\,|\,\texttt{0x0F}=\boxed{\texttt{0xFBCF}}$.
\end{enumerate}
\textit{(Náčrt doplněn k~výkladu okruhu; otázka nemá v~PDF řešení, výpočet ověřen skriptem \texttt{verify\_szz.py}.)}
